% Source: https://artofproblemsolving.com/community/c3h2451358_mathcounts_2012_state_target8
% Stable identity: aops-topic-2451358-q01
\textbf{The probability.} The $6^4=1296$ ordered outcomes are equally likely. Exactly three different numbers means that one number appears twice and two other numbers appear once each. Choose the repeated number ($6$ ways) and the two dice that show it ($\dbinom42=6$ ways); the other two dice, in order, show two different numbers from the remaining five ($5\cdot4=20$ ways). So there are $6\cdot6\cdot20=720$ favourable outcomes, and the probability is
\[\dfrac{720}{1296}=\boxed{\dfrac59}.\]

\textbf{The error.} The student's product is the probability of one particular pattern: the first three dice all different and the fourth die repeating one of them. But the repeated number can also appear among the first three dice, with the fourth die showing a new number. That pattern is missing. Its probability is found in the same way: the first three dice show exactly two different numbers with probability
\[1-P(\text{all three equal})-P(\text{all three different})=1-\dfrac{1}{36}-\dfrac56\cdot\dfrac46=1-\dfrac{1}{36}-\dfrac{20}{36}=\dfrac{15}{36}=\dfrac{5}{12},\]
and then the fourth die must differ from both numbers, with probability $\dfrac46=\dfrac23$; together $\dfrac{5}{12}\cdot\dfrac23=\dfrac{5}{18}$. The two patterns do not overlap and together cover every outcome with exactly three different numbers, so the probability is $\dfrac{5}{18}+\dfrac{5}{18}=\dfrac59$.

In terms of positions: the two equal dice can be any of the $\dbinom42=6$ pairs of dice, but the student's product only counts the $3$ pairs that include the fourth die, which is why it is exactly half of the correct answer.
